\section{Supplementary: Per-Taxon Analysis}

\label{app:per_taxon}
% ══════════════════════════════════════════════════════════════════════════════

\begin{table}[H]
\centering
\caption{Detailed per-taxon performance breakdown (5-way 5-shot \%).}
\label{tab:per_taxon}
\begin{tabular}{llcc}
\toprule
\textbf{Dataset} & \textbf{Difficulty} & \textbf{TaxoNet} &
  \textbf{P$>$M$>$F} \\
\midrule
\multirow{3}{*}{ZooBirds}
  & Inter-family & 97.8 & 96.1 \\
  & Inter-genus & 93.6 & 88.4 \\
  & Intra-genus & 86.2 & 78.9 \\
\midrule
\multirow{3}{*}{ZooInsects}
  & Inter-family & 96.4 & 94.2 \\
  & Inter-genus & 90.1 & 83.7 \\
  & Intra-genus & 81.3 & 71.6 \\
\midrule
\multirow{3}{*}{ZooPlants}
  & Inter-family & 97.1 & 95.3 \\
  & Inter-genus & 92.8 & 86.1 \\
  & Intra-genus & 84.7 & 74.8 \\
\midrule
\multirow{3}{*}{ZooMarine}
  & Inter-family & 97.3 & 95.7 \\
  & Inter-genus & 92.4 & 87.0 \\
  & Intra-genus & 85.1 & 76.2 \\
\bottomrule
\end{tabular}
\end{table}

TaxoNet's advantage is most pronounced for intra-genus discrimination
(+7.3\% to +9.9\% over P$>$M$>$F), precisely the setting where
taxonomic prototypes provide the most value.

